\subsection{Matrix Factorization and Topic Models}

Many unsupervised models express data as combinations of latent factors.

\begin{definition} \label{def:sl-low-rank-matrix-factorization} \hlt{Low-Rank Matrix Factorization}\\
For a data matrix $X\in\R^{m\times n}$, a rank-$r$ factorization approximates
\[
X\approx UV^\top,
\qquad
U\in\R^{m\times r},\quad V\in\R^{n\times r}.
\]
The factors may represent latent coordinates, user-item preferences, or basis dictionaries.
\end{definition}

\begin{definition} \label{def:sl-nonnegative-matrix-factorization} \hlt{Nonnegative Matrix Factorization}\\
Nonnegative matrix factorization imposes $U,V\geq0$ entrywise.  The constraint encourages additive parts-based representations when the data are counts or intensities.
\end{definition}

\begin{definition} \label{def:sl-topic-model} \hlt{Topic Model}\\
A topic model represents a document as a mixture of latent topics, with each topic a distribution over words.  Latent Dirichlet allocation places Dirichlet priors on document-topic and topic-word distributions.
\end{definition}

\begin{remark} \label{rem:sl-factorization-regularization} \hlt{Regularized Factors}\\
Low-rank, sparse, nonnegative, and probabilistic constraints turn an underdetermined factorization into a learning problem with interpretable inductive bias.
\end{remark}
